Let $A$ be a finite multiset of (typically positive) integers, and write
\[
R(A):=\sum_{n\in A}\frac1n.
\]
For a submultiset $S\subseteq A$ write $R(S)=\sum_{n\in S}1/n$. Define the ``gap below $1$'' as
\[
\varepsilon(A):=\min\{\,1-R(S):\ S\subseteq A,\ R(S)\le 1\,\}.
\]
(Thus $0\le \varepsilon(A)\le 1$, and $\varepsilon(A)=0$ iff some submultiset sums to $1$ exactly.)

The question asks whether there exists an absolute constant $c>0$ such that for all $K>1$,
whenever $A$ is a sufficiently large finite multiset with $R(A)>K$, one has
\[
\varepsilon(A)\le e^{-cK},
\]
equivalently: there exists $S\subseteq A$ with
\[
1-e^{-cK}<R(S)\le 1.
\]

\textbf{Stress points / ambiguity checks.}
\begin{itemize}
\item If $1\in A$ then $\varepsilon(A)=0$ trivially by taking $S=\{1\}$; the interesting regime is $A\subseteq\{2,3,4,\dots\}$.
\item The phrase ``sufficiently large'' is ambiguous in isolation (large cardinality? large elements?); I interpret it in the standard way: there exists $N_0(K)$ such that the conclusion holds whenever $|A|\ge N_0(K)$ and $R(A)>K$.
\item The problem is quantitative: Erd\H{o}s--Graham reportedly knew a weaker bound with $e^{-cK}$ replaced by $c/K^2$.
\end{itemize}
